%% ma: L12-B08-0017
%% lop: 12
%% chuong: 2
%% bai: 8
%% dang: 12.8.12
%% loai: TN4
%% muc: TH
%% dap_an: "A"
%% nguon: "(NBV)-ĐỀ SỐ 10-MỖI NGÀY 1 ĐỀ THI 2025.pdf (D10, MỖI NGÀY 1 ĐỀ - Đề số 10), Phần I Câu 11"
%% kien_thuc: "Bài 8 (tọa độ vectơ, hai vectơ bằng nhau), Bài 7 (tọa độ điểm)"
%% trang_thai: cho-duyet
%% ly_do: "Dạng mới đề xuất, chờ giáo viên duyệt."
\begin{ex}
Trong không gian $Oxyz$, cho hình bình hành $ABCD$ có $A(-3;1;2)$, $B(-2;4;-1)$, $C(1;-3;3)$. Tọa độ điểm $D$ là
\choice{\True $(0;-6;6)$}{$(-6;7;-2)$}{$(2;0;0)$}{$(-4;2;5)$}
\loigiai{$ABCD$ là hình bình hành nên $\overrightarrow{AD}=\overrightarrow{BC}$. Ta có $\overrightarrow{BC}=(3;-7;4)$. Gọi $D(x;y;z)$ thì $(x+3;y-1;z-2)=(3;-7;4)$, suy ra $D(0;-6;6)$.}
\end{ex}
